\documentclass[10pt,a4paper]{amsart}
\usepackage[margin=19mm]{geometry}
\usepackage{amsmath,amssymb,mathtools}
\usepackage[scr=rsfs,cal=euler]{mathalfa}
\usepackage{xfrac}
\usepackage[unicode,hidelinks,bookmarks=false]{hyperref}
\newcommand{\ie}{i.e.\ }
\newcommand{\cf}{cf.\ }
\newcommand{\resp}{respectively\ }
\newcommand{\Subsection}{\subsection}
\providecommand{\nobreakdash}{\nobreak\hbox{-}}
\setlength{\emergencystretch}{2em}
\multlinegap0pt
\usepackage{bm}
\usepackage{bbm}
\usepackage{dsfont}
\usepackage{xspace}
\usepackage{cancel}
\usepackage{mathtools}
\usepackage{amsthm}
\makeatletter
\@ifundefined{lemma}{\newtheorem{lemma}{Lemma}}{}
\makeatother
\title[Matrix bordering structure of the Faddeev-Jackiw algorithm: ]{Matrix bordering structure of the Faddeev-Jackiw algorithm: kernel reduction and symbolic automation}
\author{E. Chan-L\'opez, A. Mart\'in-Ruiz, Jaime Manuel Cabrera, Jorge Mauricio Paulin Fuentes}
\date{Source: arXiv:2602.12114v3 (CC BY 4.0)}
\begin{document}
\maketitle

\noindent\emph{Source and licence.} Matrix bordering structure of the Faddeev-Jackiw algorithm: kernel reduction and symbolic automation. Version arXiv:2602.12114v3 (CC BY 4.0), distributed under the Creative Commons Attribution 4.0 International licence, as recorded in the arXiv licence metadata for this exact version and shown on the abstract page https://arxiv.org/abs/2602.12114v3. Full rights notice: licenses/arxiv-2602.12114-CC-BY-4.0.txt.\par\smallskip

\noindent\textit{Editorial scope.} Counted unit: the Lemma whose source label is \texttt{lem:rank} in the article (source0.tex lines 948--967, source label \texttt{lem:rank}), delivered together with the article's own complete original proof of it (source0.tex lines 969--990, one \texttt{proof} environment of 22 lines). No mathematics is written, repaired or re-derived: statement and proof are the article's own text line for line. Required context retained uncounted: none. Every definition, display and label of those lines is retained unchanged. Omitted from this package and not redistributed: 1--947, 968--968, 991--1207. Nothing inside a retained excerpt is omitted or altered: every retained body is delivered exactly as the article's lines are written, so no omission receipt of any kind accompanies this package. Citations: the retained excerpts carry no citation command at all, so no bibliography entry of the article is transcribed and no reference is redistributed. Reference adaptation: none was needed and none was made; every reference inside the delivered unit resolves inside it, so no unresolved reference or citation is produced. Printed slips kept literal: no printed word, symbol, hypothesis or number of the counted statement or of its proof was altered. Layout and class: the article's own class is \texttt{article} with packages including inputenc, fontenc, babel, amsmath, amssymb, amsthm, mathtools, physics, graphicx, float, booktabs, caption, microtype, hyperref; the wrapper is the lane's pinned \texttt{amsart} editorial wrapper at 10pt on A4, which restates the theorem-like environments the retained text needs and adds the article's own macro definitions that the retained text needs (none), with extra wrapper packages (bm, bbm, dsfont, xspace, cancel, mathtools). The delivered numbering is the unit's own. Assets: no retained span contains a figure environment, a graphics-inclusion command, a PDF-page inclusion, a table or a TikZ picture; no image file is copied into this package, the package declares no asset dependency and no image file is redistributed with it. Licence: version arXiv:2602.12114v3 of this article is distributed under the Creative Commons Attribution 4.0 International licence, as recorded in the arXiv licence metadata for this exact version and shown on the abstract page https://arxiv.org/abs/2602.12114v3. A full rights notice is delivered at licenses/arxiv-2602.12114-CC-BY-4.0.txt. Characters: every retained excerpt is pure ASCII, so the delivered engine, XeLaTeX with its default Latin Modern text font, reports no missing character.
\begin{lemma}[Determinantal factorization for antisymmetric bordered matrices]\label{lem:rank}
	Let $f^{(0)} \in \mathbb{R}^{n \times n}$ be real antisymmetric with $\mathrm{rank}(f^{(0)}) = r = n - d$, $d = \dim\ker(f^{(0)}) \geq 1$, and let $B \in \mathbb{R}^{n \times k}$. Let $N \in \mathbb{R}^{n\times d}$ have orthonormal columns spanning $\ker(f^{(0)})$, and define the \emph{reduced constraint matrix}
	\[
	\Gamma \;\equiv\; N^{\top} B \;\in\; \mathbb{R}^{d\times k}.
	\]
	Then the extended symplectic matrix
	\[
	f^{(m)} \;=\; \begin{pmatrix} f^{(0)} & B \\ - B^{\top} & 0 \end{pmatrix}
	\]
	satisfies $\det(f^{(m)}) = \det(J)\,\det(R)$, where $J$ is the nonsingular block of $f^{(0)}$ and $R=\begin{psmallmatrix} 0 & \Gamma \\ -\Gamma^{\top} & M\end{psmallmatrix}$ with $M$ antisymmetric. In particular, if $k = d$ then
	\begin{equation}\label{eq:rank-lemma}
		\det\bigl(f^{(m)}\bigr) \;=\; \Bigl(\textstyle\prod_{\ell=1}^{r/2}\mu_\ell^{2}\Bigr)\,\det(\Gamma)^2 ,
	\end{equation}
	where $\mu_\ell>0$ are the Pfaffian factors of $f^{(0)}$. Consequently $f^{(m)}$ is nonsingular (i.e.\ of rank $n+k$) if and only if:
	\begin{enumerate}
		\item[\textup{(i)}] $k = d = n - r$, and
		\item[\textup{(ii)}] $\det(\Gamma) \neq 0$, i.e.\ $\mathrm{rank}(\Gamma) = d$,
	\end{enumerate}
	that is, if and only if the number of independent generated constraints equals $\dim\ker(f^{(0)})$ and their gradients pair non-degenerately with the null space of $f^{(0)}$.
\end{lemma}

\begin{proof}
	Since $f^{(0)}$ is real antisymmetric of rank $r$, there exists an orthogonal $Q \in \mathbb{R}^{n\times n}$ such that
	\[
	Q^{\top} f^{(0)} Q \;=\; \begin{pmatrix} J & 0 \\ 0 & 0 \end{pmatrix}, \qquad J = \bigoplus_{\ell=1}^{r/2}\mu_\ell\!\begin{psmallmatrix}0 & 1 \\ -1 & 0\end{psmallmatrix}\ \text{invertible},\quad \mu_\ell>0,
	\]
	and we may take the last $d$ columns of $Q$ to be the orthonormal kernel basis $N$. Partition $Q^{\top} B = (B_u^{\top}, B_w^{\top})^{\top}$ with $B_u \in \mathbb{R}^{r\times k}$ and $B_w \in \mathbb{R}^{d\times k}$; by construction $B_w = N^{\top}B = \Gamma$.
	
	Conjugating $f^{(m)}$ by $\mathrm{diag}(Q, I_k)$ yields
	\[
	\widetilde{f}^{(m)} \;=\; \begin{pmatrix} J & 0 & B_u \\ 0 & 0 & \Gamma \\ -B_u^{\top} & -\Gamma^{\top} & 0 \end{pmatrix}.
	\]
	Because $J$ is invertible, the congruence with the unit upper-triangular
	$T = \begin{psmallmatrix} I_r & 0 & -J^{-1}B_u \\ 0 & I_{d} & 0 \\ 0 & 0 & I_k \end{psmallmatrix}$
	(with $\det T = 1$) eliminates the $(1,3)$ and $(3,1)$ blocks:
	\[
	T^{\top}\, \widetilde{f}^{(m)}\, T \;=\; \begin{pmatrix} J & 0 & 0 \\ 0 & 0 & \Gamma \\ 0 & -\Gamma^{\top} & M \end{pmatrix},
	\qquad M = B_u^{\top} J^{-1} B_u\ \text{(antisymmetric)}.
	\]
	Since congruence preserves the determinant up to $(\det T)^2 = 1$, we obtain $\det(f^{(m)}) = \det(J)\,\det(R)$ with $R = \begin{psmallmatrix} 0 & \Gamma \\ -\Gamma^{\top} & M \end{psmallmatrix}$ and $\det(J) = \prod_{\ell}\mu_\ell^{2}$.
	
	It remains to evaluate $\det(R)$ when $k=d$. Then $R$ is antisymmetric of even size $2d$, so $\det(R)=\mathrm{Pf}(R)^2$. The block congruence with $\begin{psmallmatrix} I_d & 0 \\ X & I_d \end{psmallmatrix}$, where $X$ is chosen so that $X\Gamma - (X\Gamma)^{\top} = M$ (solvable because $M$ is antisymmetric and $\Gamma$ invertible), maps $R$ to $\begin{psmallmatrix} 0 & \Gamma \\ -\Gamma^{\top} & 0 \end{psmallmatrix}$ without changing its Pfaffian; hence $\mathrm{Pf}(R) = (-1)^{d(d-1)/2}\det(\Gamma)$ and $\det(R)=\det(\Gamma)^2$, independently of the image block $M$. This proves~\eqref{eq:rank-lemma}. Since $\prod_\ell\mu_\ell^2>0$, $f^{(m)}$ is nonsingular iff $\det(\Gamma)\neq 0$; and $\Gamma\in\mathbb{R}^{d\times k}$ can have $\det\neq 0$ only when it is square, i.e.\ $k=d$, giving conditions (i) and (ii). If $k\neq d$ the size $n+k$ has the parity of $n-r+k$ and $\Gamma$ is not square, so $f^{(m)}$ is singular.
\end{proof}

\end{document}
